% Source: https://artofproblemsolving.com/community/c3h2977061_math_problem_set_2
% Stable identity: aops-topic-2977061-q01
Let $x=\dfrac{a(a+1)}2$, $y=\dfrac{b(b+1)}2$ with $a>b\ge1$. Then $x-y=\dfrac{(a-b)(a+b+1)}2=14$, so $(a-b)(a+b+1)=28$. The two factors have opposite parity, and $a+b+1>a-b$. The options are $1\cdot28$ and $4\cdot7$:
\begin{itemize}
\item $a-b=1$, $a+b=27$: $a=14$, $b=13$, $x=105$, $y=91$, $z=392$.
\item $a-b=4$, $a+b=6$: $a=5$, $b=1$, $x=15$, $y=1$, $z=32$.
\end{itemize}
So $z=\boxed{32\text{ or }392}$.
